\documentclass[aspectratio=169, 10pt]{beamer}

% --- Typography and Encoding ---
\usepackage[utf8]{inputenc}
\usepackage[T1]{fontenc}
\usepackage{lmodern}
\usepackage{microtype}

% --- Math & Tables ---
\usepackage{amsmath, amssymb, amsfonts}
\usepackage{booktabs}
\usepackage{tabularx}
\usepackage{array}

% --- Graphics & TikZ ---
\usepackage{graphicx}
\usepackage{tikz}
\usetikzlibrary{shapes.geometric, arrows.meta, positioning, calc}

% --- Color Palette (Keio Navy, Slate, Academic Accents) ---
\definecolor{keioblue}{RGB}{0, 32, 91}       % #00205B
\definecolor{darkslate}{RGB}{30, 41, 59}     % #1E293B
\definecolor{accentred}{RGB}{185, 28, 28}    % #B91C1C
\definecolor{softgrey}{RGB}{241, 245, 249}   % #F1F5F9
\definecolor{linegrey}{RGB}{203, 213, 225}   % #CBD5E1
\definecolor{textgrey}{RGB}{100, 116, 139}   % #64748B

% --- Beamer Theme Customization ---
\setbeamercolor{title}{fg=keioblue}
\setbeamercolor{frametitle}{fg=keioblue, bg=white}
\setbeamercolor{structure}{fg=keioblue}
\setbeamercolor{normal text}{fg=darkslate}
\setbeamercolor{block title}{fg=white, bg=keioblue}
\setbeamercolor{block body}{fg=darkslate, bg=softgrey}
\setbeamercolor{item}{fg=keioblue}

\setbeamertemplate{navigation symbols}{}
\setbeamertemplate{itemize items}[circle]

% --- Clean Header Rule ---
\setbeamertemplate{frametitle}{%
    \vspace*{0.3cm}
    \bfseries\insertframetitle
    \par\vspace*{0.1cm}
    {\color{linegrey}\hrule height 0.5pt}
}

% --- Custom Footline (Current / Total strictly for main deck) ---
\setbeamertemplate{footline}{%
    \leavevmode%
    \hbox{%
        \begin{beamercolorbox}[wd=0.32\paperwidth, ht=2.5ex, dp=1.1ex, left]{}%
            \hspace*{1.5em}\scriptsize\textcolor{textgrey}{Eliott Flechtner (M2 QMI)}
        \end{beamercolorbox}%
        \begin{beamercolorbox}[wd=0.48\paperwidth, ht=2.5ex, dp=1.1ex, center]{}%
            \scriptsize\textcolor{textgrey}{Purification Scheduling for All-Photonic Repeaters}
        \end{beamercolorbox}%
        \begin{beamercolorbox}[wd=0.20\paperwidth, ht=2.5ex, dp=1.1ex, right]{}%
            \scriptsize\textcolor{textgrey}{\insertframenumber{} / \inserttotalframenumber}\hspace*{1.5em}
        \end{beamercolorbox}%
    }%
    \vskip0pt%
}

% --- Macro for Standard Graphic Placeholders ---
\newcommand{\graphicplaceholder}[3]{%
    \begin{tikzpicture}
        \node[draw=linegrey, dashed, fill=softgrey, rounded corners=2pt, minimum width=#1, minimum height=#2, align=center, inner sep=6pt] {
            \textcolor{textgrey}{\small\textbf{#3}}\\[2pt]
            \textcolor{textgrey}{\scriptsize[Replace with asset / PDF figure]}
        };
    \end{tikzpicture}
}

% ==============================================================================
% PRESENTATION METADATA
% ==============================================================================
\title{Non-Adaptive Scheduling of Entanglement Purification\\for All-Photonic Quantum Repeaters}
\subtitle{Master's Thesis Defense --- M2 QMI}
\author{Eliott Flechtner}
\institute{EPITA \& Institut Polytechnique de Paris\\Host Laboratory: AQUA, Keio University SFC Campus, Japan}
\date{September 25, 2026}

\begin{document}

% ==============================================================================
% SLIDE 1: TITLE SLIDE
% ==============================================================================
{
\setbeamertemplate{footline}{}
\begin{frame}
    % Institutional Header Logos Placeholder
    \vspace*{-0.2cm}
    \begin{columns}[c]
        \column{0.30\textwidth}
            \centering \graphicplaceholder{3.2cm}{0.9cm}{Keio University}
        \column{0.40\textwidth}
            \centering \graphicplaceholder{4.2cm}{0.9cm}{Télécom Paris / IP Paris}
        \column{0.30\textwidth}
            \centering \graphicplaceholder{3.2cm}{0.9cm}{EPITA}
    \end{columns}

    \vspace{0.6cm}
    \centering
    {\large\bfseries\textcolor{keioblue}{Non-Adaptive Scheduling of Entanglement Purification\\for All-Photonic Quantum Repeaters}}\\[0.25cm]
    {\footnotesize\textcolor{darkslate}{Research Internship Defense --- Master 2 Quantum, Mathematics and Informatics (QMI)}}\\[0.15cm]
    {\normalsize\textbf{Eliott Flechtner}}\\[0.1cm]
    {\scriptsize EPITA \& Institut Polytechnique de Paris \quad$\vert$\quad AQUA Lab, Keio University, Japan}

    \vspace{0.4cm}
    \begin{columns}[t]
        \column{0.50\textwidth}
            \raggedright
            \tiny\textbf{\textcolor{textgrey}{SUPERVISORY COMMITTEE}}\\[2pt]
            \scriptsize
            Prof. Rodney Van Meter \textit{(Keio University)}\\
            Proj. Assoc. Prof. Michal Hajdušek \textit{(Keio University)}\\
            Proj. Assist. Prof. Naphan Benchasattabuse \textit{(Keio University)}
        \column{0.45\textwidth}
            \raggedright
            \tiny\textbf{\textcolor{textgrey}{EXAMINATION COMMITTEE}}\\[2pt]
            \scriptsize
            Prof. Romain Alléaume \textit{(Télécom Paris, IP Paris)}\\
            Prof. Daniel E. Browne \textit{(University College London)}
    \end{columns}
\end{frame}
}

% ==============================================================================
% SLIDE 2: CONTEXT & MOTIVATION
% ==============================================================================
\begin{frame}{Why Quantum Repeaters}
    \begin{columns}[T]
        \column{0.55\textwidth}
            \begin{itemize}
                \item Fiber loss grows exponentially with distance
                \item Direct transmission cannot scale
                \item Repeaters break the channel into shorter, purified hops
            \end{itemize}
            \vspace{0.2cm}
            \begin{itemize}
                \item 1G/2G: memory-based, limited by dephasing \\\& two-way signaling
                \item 3G, all-photonic: flying graph states, loss tolerance via tree encoding
            \end{itemize}
        \column{0.42\textwidth}
            \centering
            \graphicplaceholder{5.5cm}{4.5cm}{Fiber loss vs.\ all-photonic repeater concept}
    \end{columns}
\end{frame}

% ==============================================================================
% SLIDE 3: REPEATER ARCHITECTURE
% ==============================================================================
\begin{frame}{Half-RGS Repeaters}
    \begin{columns}[T]
        \column{0.55\textwidth}
            \begin{itemize}
                \item HRGS: half of a repeater graph state, held at one station
                \item End nodes act as their own entanglement source
                \item Memory cost grows additively with hop count, not multiplicatively
            \end{itemize}
            \vspace{0.2cm}
            \begin{itemize}
                \item Outer photons meet at analyzer stations
                \item Linear-optical BSM success capped at 50\% $\to$ redundant arms
            \end{itemize}
        \column{0.42\textwidth}
            \centering
            \graphicplaceholder{5.5cm}{4.5cm}{HRGS combination and $N$-hop chain}
    \end{columns}
\end{frame}

% ==============================================================================
% SLIDE 4: THE SCHEDULING GAP
% ==============================================================================
\begin{frame}{The Scheduling Gap}
    \begin{columns}[T]
        \column{0.58\textwidth}
            \begin{itemize}
                \item Purification can be inserted at generation, link, intermediate span, or end-to-end
                \item Circuit, copy count and order can vary between stages
            \end{itemize}
            \vspace{0.3cm}
            \begin{itemize}
                \item Prior work (Integrating paper) shows this flexibility with one hand-picked schedule
                \item Which allocation is actually best under a fixed budget was left open
            \end{itemize}
        \column{0.38\textwidth}
            \centering
            \graphicplaceholder{5.0cm}{4.5cm}{Purification insertion points across spans}
    \end{columns}
\end{frame}

% ==============================================================================
% SLIDE 5: ERROR REPRESENTATION & CIRCUITS
% ==============================================================================
\begin{frame}{Noise Model and Purification Circuits}
    \begin{columns}[T]
        \column{0.50\textwidth}
            \begin{itemize}
                \item Bell-diagonal states, error vector $e=(w,x,y,z)$, fidelity $F=w$
                \item Inner-qubit errors bias $x,z$ but leave $y\approx0$
            \end{itemize}
        \column{0.48\textwidth}
            \begin{itemize}
                \item Three 2-to-1 circuits: $C_{ZX}$, $C_{XZ}$, $C_{YY}$
                \item $C_{YY}$ used first to remove the single-sided bias
            \end{itemize}
            \vspace{0.2cm}
            \centering
            \graphicplaceholder{5.0cm}{1.5cm}{Circuit layouts}
    \end{columns}
\end{frame}

% ==============================================================================
% SLIDE 6: PARTIALLY-BUILT BELL STATE S
% ==============================================================================
\begin{frame}{Representing a Partial Resource}
    Every intermediate resource is a state:
    $$S = \big(e,\; s,\; t,\; t_{\text{gen}},\; \kappa,\; r,\; h\big)$$

    \begin{columns}[T]
        \column{0.48\textwidth}
            \begin{itemize}
                \item $e$: error vector, $s$: Pauli frame bit
                \item $t-t_{\text{gen}}$: idle time $\to$ memory dephasing
            \end{itemize}
        \column{0.48\textwidth}
            \begin{itemize}
                \item $\kappa$: span, $r$: purification rounds, $h$: heralding status
            \end{itemize}
            \vspace{0.2cm}
            \begin{alertblock}{Span-matching rule}
                Two states purify together only if $\kappa(S_1)=\kappa(S_2)$ --- construction history doesn't matter.
            \end{alertblock}
    \end{columns}
\end{frame}

% ==============================================================================
% SLIDE 7: SCHEDULES AS DAGS
% ==============================================================================
\begin{frame}{A Schedule Is a DAG}
    \begin{columns}[T]
        \column{0.45\textwidth}
            $\Sigma = (T, \phi)$, $\phi: T \to \mathcal{K}$
            \vspace{0.2cm}
            \begin{itemize}
                \item Gen, Join, Swap, Purify, Idle, Herald, PauliCorrect
            \end{itemize}
            \vspace{0.2cm}
            \begin{itemize}
                \item Optimistic: heralding deferred to a single node at the root
            \end{itemize}
        \column{0.52\textwidth}
            \centering
            \graphicplaceholder{6.5cm}{4.6cm}{$N=2$ worked example DAG}
    \end{columns}
\end{frame}

% ==============================================================================
% SLIDE 8: CONCURRENCY & EVALUATION PASS
% ==============================================================================
\begin{frame}{Evaluating a Schedule}
    \begin{columns}[T]
        \column{0.50\textwidth}
            \begin{itemize}
                \item Single bottom-up pass, $O(|T|)$: fidelity, cost, latency
                \item Rate accounts for full-schedule restart on failure:
                $$R(\Sigma) = \frac{P_{\text{success}}}{T_{\text{attempt}}}$$
            \end{itemize}
            \vspace{0.2cm}
            $$\Sigma^* = \arg\max_{\Sigma \text{ feasible}} R(\Sigma) \;\; \text{s.t.}\;\; F(\Sigma) \ge F_{\min}$$
        \column{0.48\textwidth}
            \begin{itemize}
                \item Concurrency $M(\Sigma)$: min. open branches, via Sethi--Ullman recurrence
                \item Exact register bound, no timing info needed
            \end{itemize}
    \end{columns}
\end{frame}

% ==============================================================================
% SLIDE 9: COMBINATORIAL SEARCH FRAMEWORK
% ==============================================================================
\begin{frame}{Searching the Schedule Space}
    \begin{columns}[T]
        \column{0.55\textwidth}
            \begin{itemize}
                \item Combinatorial explosion: circuit choice, split points, copy counts
                \item Exhaustive search intractable beyond $N=3$
            \end{itemize}
            \vspace{0.3cm}
            \textbf{Three tiers:}
            \begin{enumerate}
                \item Fixed baselines (raw, uniform link-level, paper schedule)
                \item Exact span-partition dynamic program
                \item Bounded beam search with same-span pumping
            \end{enumerate}
        \column{0.42\textwidth}
            \centering
            \graphicplaceholder{5.5cm}{4.5cm}{Baselines $\to$ exact DP $\to$ beam search}
    \end{columns}
\end{frame}

% ==============================================================================
% SLIDE 10: TIER 2: SPAN DYNAMIC PROGRAMMING
% ==============================================================================
\begin{frame}{Exact Search: Span-Partition DP}
    \begin{columns}[T]
        \column{0.52\textwidth}
            \begin{itemize}
                \item Memoize best sub-schedules per span, reuse to build wider spans:
                $$\text{Span}(a, c) = \text{Span}(a, b) \bowtie \text{Span}(b, c)$$
            \end{itemize}
            \vspace{0.2cm}
            \begin{itemize}
                \item Keep a Pareto set over (cost, fidelity, success prob.)
                \item Even a locally worse candidate can win later once swapped
            \end{itemize}
        \column{0.45\textwidth}
            \centering
            \graphicplaceholder{5.5cm}{4.2cm}{Span memoization \& Pareto pruning}
    \end{columns}
\end{frame}

% ==============================================================================
% SLIDE 11: TIER 3: SAME-SPAN PUMPING & BEAM SEARCH
% ==============================================================================
\begin{frame}{Pumping and Beam Search}
    \begin{columns}[T]
        \column{0.55\textwidth}
            \begin{itemize}
                \item Same-span pumping: purify against a fresh copy of the same span, not just a different one
            \end{itemize}
            \vspace{0.2cm}
            \begin{itemize}
                \item Beam width 25 keeps frontiers tractable (uncapped: 1000+ candidates)
                \item Half the slots for top fidelity, half for best rate among feasible candidates
            \end{itemize}
        \column{0.42\textwidth}
            \centering
            \graphicplaceholder{5.5cm}{4.5cm}{Beam frontier retention strategy}
    \end{columns}
\end{frame}

% ==============================================================================
% SLIDE 12: PHYSICS VALIDATION
% ==============================================================================
\begin{frame}{Validating the Evaluator}
    \begin{columns}[T]
        \column{0.50\textwidth}
            \begin{itemize}
                \item Reproduces the Integrating paper's Fig.\,5 fidelity curves
                \item Raw chain, heralded pumping, optimistic pumping
                \item Match to four significant figures
            \end{itemize}
        \column{0.48\textwidth}
            \centering
            \small
            \begin{tabular}{lccc}
                \toprule
                \textbf{Schedule} & \textbf{Ref.} & \textbf{Ours} \\
                \midrule
                Raw & $0.823$ & $\mathbf{0.8234}$ \\
                Heralded & $0.917$ & $\mathbf{0.9168}$ \\
                Optimistic & $0.929$ & $\mathbf{0.9295}$ \\
                \bottomrule
            \end{tabular}
            \vspace{0.4cm}
            \graphicplaceholder{5.2cm}{2.5cm}{Fidelity reproduction curve}
    \end{columns}
\end{frame}

% ==============================================================================
% SLIDE 13: HEADLINE RESULTS
% ==============================================================================
\begin{frame}{Beating the Paper's Schedule}
    \textbf{Reference config:} $N = 10$, $e_d = 0.01$, $F_{\min} = 0.90$, $E_{\max} = 100$

    \vspace{0.2cm}
    \centering
    \small
    \begin{tabularx}{0.95\textwidth}{lccccX}
        \toprule
        \textbf{Schedule} & \textbf{$C$} & \textbf{$F$} & \textbf{Rate} & \textbf{Rel.} & \\
        \midrule
        Paper schedule & 100 & 0.93 & 4,055.92 & 1.00$\times$ & hand-picked \\
        Matched cost & 100 & 0.92 & 4,158.14 & \textbf{+2.5\%} & same budget \\
        Budget-relaxed & \textbf{60} & 0.91 & \textbf{6,195.95} & \textbf{+52.8\%} & 60\% of cost \\
        \bottomrule
    \end{tabularx}

    \vspace{0.4cm}
    \begin{itemize}
        \item The paper schedule over-purifies everywhere uniformly
        \item Search spends the budget where it's actually needed
    \end{itemize}
\end{frame}

% ==============================================================================
% SLIDE 14: FEASIBILITY ADVANTAGE
% ==============================================================================
\begin{frame}{The Real Finding: Feasibility, Not Just Speed}
    \begin{columns}[T]
        \column{0.50\textwidth}
            \begin{itemize}
                \item Turn on inner-qubit error ($p_{\text{in}}=3\times10^{-3}$)
                \item Paper schedule drops below $F_{\min}$: $F=0.67$ (18 arms), $F=0.87$ (6 arms)
                \item Search finds a feasible schedule at the \textbf{same cost}
            \end{itemize}
            \vspace{0.3cm}
            \begin{itemize}
                \item Same pattern on 100 randomized heterogeneous networks: mean +20\% rate, never worse
                \item One weak hop with 12$\times$ noise: uniform purification fails, search concentrates effort there
            \end{itemize}
        \column{0.48\textwidth}
            \centering
            \graphicplaceholder{5.5cm}{4.6cm}{Fixed schedule vs.\ feasible searched schedule}
    \end{columns}
\end{frame}

% ==============================================================================
% SLIDE 15: CHAIN-LENGTH SCALING & TIMING
% ==============================================================================
\begin{frame}{Scaling with Chain Length and Timing}
    \begin{columns}[T]
        \column{0.50\textwidth}
            \begin{itemize}
                \item Minimum budget to hit $F\ge0.9$, for $N=10$ to $28$
                \item Fit: $E_{\max}^{\min} \approx 1.19\,N^{1.67}$ --- superlinear
                \item Paper's linear $10N$ rule breaks around $N=28$
            \end{itemize}
            \vspace{0.3cm}
            \begin{itemize}
                \item Heralded pumping: waiting copies decohere, $F$ drops $0.92\to0.25$
                \item Optimistic schedules avoid the wait, stay flat
            \end{itemize}
        \column{0.48\textwidth}
            \centering
            \graphicplaceholder{5.5cm}{4.6cm}{Budget scaling vs.\ dephasing decay}
    \end{columns}
\end{frame}

% ==============================================================================
% SLIDE 16: DISCUSSION: OPERATING REGIMES
% ==============================================================================
\begin{frame}{When Does Scheduling Actually Matter}
    \begin{columns}[T]
        \column{0.52\textwidth}
            \begin{itemize}
                \item Low noise ($e_d \le 0.005$): easy, everything clears the floor
                \item High noise ($e_d \ge 0.016$): nothing clears the floor
                \item In between: \textbf{scheduling is decisive}
            \end{itemize}
            \vspace{0.2cm}
            \begin{itemize}
                \item Pareto quality plateaus past $C\approx60$
            \end{itemize}
        \column{0.45\textwidth}
            \begin{block}{Scope}
                Beam search is heuristic --- results are existence proofs within the searched families, not global optima.
            \end{block}
            \vspace{0.2cm}
            \centering
            \graphicplaceholder{5.0cm}{2.0cm}{Pareto trade-off frontiers}
    \end{columns}
\end{frame}

% ==============================================================================
% SLIDE 17: CONCLUSION & OUTLOOK
% ==============================================================================
\begin{frame}{Conclusion}
    \begin{columns}[T]
        \column{0.55\textwidth}
            \begin{itemize}
                \item DAG-based schedule model + validated evaluator
                \item Three-tier search: baselines, exact DP, beam search
                \item Main result: scheduling buys \textbf{feasibility}, not just a rate bump
                \item Budget scales roughly as $N^{1.67}$
            \end{itemize}
            \vspace{0.3cm}
            \begin{itemize}
                \item Next: adaptive (MDP-based) scheduling, mesh topologies
            \end{itemize}
        \column{0.42\textwidth}
            \begin{block}{Code}
                \scriptsize\texttt{github.com/EliottFlechtner/\\entanglement-purification-scheduler}
            \end{block}
            \vspace{0.4cm}
            \centering
            \large\textbf{Thank you.}\\[0.2cm]
            \small\textit{Questions welcome.}
    \end{columns}
\end{frame}

% ==============================================================================
% APPENDIX / BACKUP SLIDES (Does not increment total slide count)
% ==============================================================================
\appendix

\begin{frame}{Backup: Sethi--Ullman Concurrency Algorithm}
    \begin{columns}[T]
        \column{0.55\textwidth}
            \begin{itemize}
                \item Bottom-up register assignment
                \item Evaluate the heavier branch first, hold result in one register
                \item Time $O(|T|)$, space $O(h)$
            \end{itemize}
        \column{0.42\textwidth}
            \centering
            \graphicplaceholder{5.5cm}{4.5cm}{RegisterRequirement(v) pseudocode}
    \end{columns}
\end{frame}

\begin{frame}{Backup: DP vs.\ Beam Search Runtime}
    \begin{columns}[T]
        \column{0.50\textwidth}
            \begin{itemize}
                \item Exact DP explodes beyond $N=10$
                \item $N=18$: DP times out (300s); beam search finds $R=1{,}214$ in 153.5s
                \item Beam search scales predictably up to $N=28$
            \end{itemize}
        \column{0.48\textwidth}
            \centering
            \graphicplaceholder{5.5cm}{4.5cm}{Runtime and score benchmarks}
    \end{columns}
\end{frame}

\end{document}